\documentclass[11pt,letterpaper]{article}

\usepackage[utf8]{inputenc}
\usepackage[T1]{fontenc}
\usepackage[margin=1in]{geometry}
\usepackage{tikz}
\usepackage{xcolor}
\usepackage{amssymb}
\usepackage{tabularx}
\usepackage{booktabs}
\usepackage{longtable}
\usepackage{enumitem}
\usepackage{hyperref}
\usepackage{titlesec}
\usepackage{fancyhdr}
\usepackage{parskip}

\definecolor{lux-red}{HTML}{8B1A1A}
\definecolor{lux-grey}{HTML}{4A4A4A}
\definecolor{lux-accent}{HTML}{1F3A5F}

\hypersetup{
  colorlinks=true,
  linkcolor=lux-red,
  citecolor=lux-grey,
  urlcolor=lux-red,
  pdftitle={Lux/Liquidity --- Overview of IP Licensing and Liquidity.io Fraud},
  pdfauthor={Woo Bin --- CTO, Lux Industries Inc.; VP Engineering, Satschel Inc.},
  pdfsubject={Confidential memorandum to the Board of Directors, Lux Industries Inc.},
}

\titleformat{\section}{\normalfont\Large\bfseries\color{lux-grey}}{\thesection}{0.6em}{}
\titleformat{\subsection}{\normalfont\large\bfseries\color{lux-grey}}{\thesubsection}{0.6em}{}

\pagestyle{fancy}
\fancyhf{}
\fancyhead[L]{\small\color{lux-grey}\textsc{Confidential --- Board Materials --- Attorney-Client Privileged}}
\fancyhead[R]{\small\color{lux-grey}Lux Industries Inc.}
\fancyfoot[C]{\small\color{lux-grey}\thepage}
\renewcommand{\headrulewidth}{0.2pt}

\begin{document}

% ─────────────────────────────────────────────────────────────────
% Cover page
% ─────────────────────────────────────────────────────────────────
\pagestyle{empty}
\begin{tikzpicture}[remember picture,overlay]
  \fill[black] (current page.south west) rectangle (current page.north east);

  \node[anchor=north west] at ([xshift=1.4cm,yshift=-1.4cm]current page.north west) {%
    \begin{tikzpicture}[scale=0.5]
      \fill[white] (-0.7,0.55) -- (0.7,0.55) -- (0,-0.55) -- cycle;
    \end{tikzpicture}%
  };

  \node[anchor=center,text=white,font=\fontsize{30}{38}\selectfont\bfseries,text width=14cm,align=center]
    at ([yshift=3.2cm]current page.center) {LUX INDUSTRIES};

  \node[anchor=center,text=white!70,font=\fontsize{14}{18}\selectfont,text width=14cm,align=center]
    at ([yshift=1.6cm]current page.center)
    {Overview of IP Licensing and Liquidity.io Fraud};

  \draw[white!30,line width=0.5pt]
    ([yshift=0.5cm,xshift=-5cm]current page.center) --
    ([yshift=0.5cm,xshift=5cm]current page.center);

  \node[anchor=center,text=white!40,font=\fontsize{9}{11}\selectfont]
    at ([yshift=-0.8cm]current page.center)
    {Prepared for: Cyrus Pahlavi, President \& CEO};

  \node[anchor=center,text=white!40,font=\fontsize{9}{11}\selectfont]
    at ([yshift=-1.4cm]current page.center)
    {Hunter Dupont, Managing Partner --- Strategic Transactions};

  \node[anchor=center,text=white!40,font=\fontsize{9}{11}\selectfont]
    at ([yshift=-2.0cm]current page.center)
    {Alan Donohue, General Counsel};

  \node[anchor=center,text=white!40,font=\fontsize{9}{11}\selectfont]
    at ([yshift=-2.8cm]current page.center)
    {Prepared by: Woo Bin --- CTO, Lux Industries Inc.\ /\ VP Engineering, Satschel Inc.};

  \node[anchor=center,text=white!30,font=\fontsize{9}{11}\selectfont]
    at ([yshift=-3.4cm]current page.center)
    {2026-05-25};

  \node[anchor=south west,text=white!20,font=\fontsize{7}{9}\selectfont]
    at ([xshift=1.4cm,yshift=0.8cm]current page.south west)
    {Lux Industries Inc.};

  \node[anchor=south east,text=white!20,font=\fontsize{7}{9}\selectfont]
    at ([xshift=-1.4cm,yshift=0.8cm]current page.south east)
    {CONFIDENTIAL --- ATTORNEY-CLIENT PRIVILEGED};
\end{tikzpicture}
\clearpage
\pagestyle{fancy}
\setcounter{page}{1}

% ─────────────────────────────────────────────────────────────────
\begin{flushleft}
\textbf{\textsc{Memorandum}}\\[0.5em]
\begin{tabularx}{\textwidth}{lX}
\textbf{TO:}        & Cyrus Pahlavi, President \& Chief Executive Officer, Lux Industries Inc. \\
                    & Hunter Dupont, Managing Partner --- Strategic Transactions (deal lead) \\
                    & Alan Donohue, General Counsel (legal lead) \\
\textbf{CC:}        & Board of Directors, Lux Industries Inc. \\
\textbf{FROM:}      & Woo Bin --- Chief Technology Officer, Lux Industries Inc.; concurrently Vice President of Engineering, Satschel, Inc.\ d/b/a Liquidity.io \\
\textbf{DATE:}      & 25 May 2026 \\
\textbf{RE:}        & Overview of IP Licensing and Liquidity.io Fraud --- Misappropriation of Lux foundational technology by Satschel, Inc.\ and Liquidity.io LLC; legal options, damages analysis, and proposed settlement framework \\
\textbf{PRIVILEGE:} & Confidential. Attorney-client privileged. Attorney work product. Prepared in anticipation of litigation and as a board strategy document. \\
\end{tabularx}
\end{flushleft}

\vspace{0.3em}
\begin{small}
\textit{Author's disclosure of interest:} I serve concurrently as
Chief Technology Officer of Lux Industries Inc.\ and Vice President of
Engineering at Satschel, Inc. (d/b/a Liquidity.io). This memorandum is
written from the Lux chair, in service of the Lux board's decision-
making. My Satschel role gives direct operational visibility into the
matters described; I have set them out as observed, with the
understanding that any Satschel-side response or counter-position will
need to be developed by Satschel's own counsel and that Lux's outside
litigation counsel will need to develop independent record support
for any allegations herein before filing.
\end{small}

\vspace{0.5em}\hrule height 0.6pt\vspace{1em}

% ─────────────────────────────────────────────────────────────────
\section*{I.\ Executive Summary}

\textbf{Satschel, Inc.\ and its operating affiliate Liquidity.io LLC
(collectively, ``Satschel'' or ``Liquidity'') have commercially
deployed material portions of Lux Industries Inc.'s patent-protected
and source-available proprietary technology since at least 2014, in
production, on three live blockchain networks (Liquid Devnet ChainID
8675312; Liquid Testnet ChainID 8675310; Liquid Mainnet ChainID
8675309), and in support of a registered ATS/BD/TA business that
recently raised approximately \$1.1 billion in capital at a post-money
valuation in that range --- in each case without executing the
commercial license that the governing Lux Ecosystem License v1.2
(``LEL'') and the Lux Research License with Patent Reservation v1.0
(``LRL--PR'') require.}

\textbf{Lux's investment in the misappropriated technology base
exceeds \$20 million in direct development cost recorded against Lux
books from 2014 to date, plus a multi-year contribution from senior
engineering personnel including the named author of this memorandum,
the Lux Chief Cryptographer, and others.}

\textbf{This memorandum sets out the legal record, identifies the
causes of action available to Lux against Satschel and its officers,
quantifies damages at small / medium / large scenarios, projects the
forward enterprise value of the equity component at user-scale
breakpoints of 100K, 1M, 10M, and 100M users, analyzes the
``wait vs.\ settle'' time-value question, and offers a settlement
triangulation matrix the deal lead (Mr.\ Dupont) and General Counsel
(Mr.\ Donohue) can use to negotiate.}

This is not a single-proposal document. It presents the data, the
damages range, the partnership-upside range, and the negotiating
floors --- the Lux deal team can triangulate to a specific outcome.
The final approval authority on any executed Settlement and Cross-
License Agreement ($\S$XII) rests with the President \& CEO, Mr.\
Pahlavi.

% ─────────────────────────────────────────────────────────────────
\section*{II.\ Parties and Dramatis Personae}

\subsection*{Lux Industries Inc.\ --- Decision Chain}

\begin{itemize}[leftmargin=1.5em,itemsep=0.2em]
\item \textbf{Cyrus Pahlavi} --- President \& Chief Executive Officer.
Final approval authority on any settlement, license, or litigation
decision.
\item \textbf{Hunter Dupont} --- Largest investor in Lux Industries
Inc.; Managing Partner, Strategic Transactions. Leads deal discussions
and commercial negotiations with Satschel.
\item \textbf{Alan Donohue} --- General Counsel. Leads legal strategy
and definitive-document drafting; provides legal support to Mr.\
Dupont throughout the negotiation.
\item \textbf{Woo Bin} --- Chief Technology Officer. Author of this
memorandum. Holds concurrent role as Vice President of Engineering at
Satschel, Inc.; operational visibility into Satschel's use of Lux IP.
\item \textbf{Zach Kelling} --- Founder, Architect, and Chief
Cryptographer of Lux. Minority shareholder. \textbf{See \S X for
particularized treatment.} Not a director on either side. Not the
current CTO of either entity.
\end{itemize}

\subsection*{Satschel, Inc.\ d/b/a Liquidity.io --- Counterparty Officer Slate}

\begin{itemize}[leftmargin=1.5em,itemsep=0.2em]
\item \textbf{Eric Choi} --- Chief Executive Officer. Top operational
authority and the officer with ultimate responsibility for Satschel's
representations to investors and counterparties regarding IP ownership
and licensing.
\item \textbf{Austin Trombley} --- Chief Strategy Officer; previously
Acting Chief Executive Officer. Blockchain / AI engineer. Built a
\textbf{prior} (not current) ATS infrastructure at a prior employer
(Franklin Templeton); served as an eyewitness in the Franklin Templeton
litigation arising from that engagement. \textbf{Counsel should
review the Franklin Templeton litigation record} for (i) pattern-of-
conduct evidence relevant to the present matter, (ii) any limits on or
adverse characterizations of his prior testimony that could be
impeachable, and (iii) the technical scope and quality of the ATS he
built there. Material context: prior to Mr.\ Kelling's arrival as
active Chief Technology Officer at Satschel (see infra \S IV(G)),
the Satschel ATS as built under Mr.\ Trombley's technical direction
\textbf{lacked even basic time-price priority order matching, the
foundational principle of any functioning alternative trading
system}. The technical state of the pre-Kelling Satschel platform is
addressed in \S IV(G) and bears directly on Mr.\ Trombley's
credibility as a technology witness in any subsequent proceeding.
\item \textbf{Coleman Church} --- Chief Executive Officer, Liquidity
Broker-Dealer. Responsible for BD-side compliance posture and FINRA
representations.
\item \textbf{Jayant Sonsurkar} --- Chief Security Officer. Long-
standing maintainer / admin on the \texttt{luxfi/*} GitHub organization.
The CSO role at any regulated digital-asset venue carries the duty to
oversee software-supply-chain integrity, including open-source license
review and third-party-IP compliance. As maintainer on every Lux
repository Satschel later deployed, Mr.\ Sonsurkar had on-platform
visibility into the LEL, LRL--PR, BSD-3, and other LICENSE files as
they were created, amended, and committed. The license tier of each
Lux module was directly observable in commit history under his admin
access.
\item \textbf{Zach Kelling} --- Founder of Satschel; Architect; Chief
Cryptographer. \textbf{Not the current CEO or CTO of Satschel;
previously held those titles in earlier corporate periods. Minority
shareholder of Lux Industries Inc.\ and \emph{zero}-holder of any
Satschel / Liquidity.io equity}. Effectively, on the Satschel side,
Mr.\ Kelling has been operationally displaced from the executive
roles he previously held and now occupies an architect / chief-
cryptographer position with limited day-to-day decision authority over
the commercial direction. \textbf{See \S X.}
\end{itemize}

\subsection*{Counterparty Entity Structure}

For purposes of this memorandum, ``Satschel'' and ``Liquidity''
refer collectively to Satschel, Inc.\ (a Delaware corporation) and
Liquidity.io LLC (its operating affiliate), unless context requires
otherwise. The corporate structure should be confirmed during outside-
counsel diligence.

% ─────────────────────────────────────────────────────────────────
\section*{III.\ Jurisdiction, Venue, and Applicable Law}

\textbf{Federal causes of action:}

\begin{itemize}[leftmargin=1.5em,itemsep=0.1em]
\item \textbf{Defend Trade Secrets Act (DTSA)} --- 18 U.S.C.\ \S\S
1836--1839. Federal question jurisdiction under 28 U.S.C.\ \S 1331;
nationwide service of process under DTSA \S 1836(c).
\item \textbf{Copyright Act} --- 17 U.S.C.\ \S\S 101 \emph{et seq.}
Exclusive federal jurisdiction under 28 U.S.C.\ \S 1338(a).
\item \textbf{Patent Act} --- 35 U.S.C.\ \S\S 271, 281, 284, 285.
Exclusive federal jurisdiction; Federal Circuit on appeal.
\item \textbf{Lanham Act} --- 15 U.S.C.\ \S 1125(a), false advertising
/ false designation of origin.
\item \textbf{Securities Exchange Act of 1934} --- 15 U.S.C.\ \S
78j(b); Rule 10b-5, 17 C.F.R.\ \S 240.10b-5. Federal question
jurisdiction.
\end{itemize}

\textbf{State / common-law causes of action:}

\begin{itemize}[leftmargin=1.5em,itemsep=0.1em]
\item \textbf{Delaware Breach of Fiduciary Duty} --- against Satschel
officers individually. Court of Chancery; see \emph{Aronson v.\
Lewis}, 473 A.2d 805 (Del.\ 1984); \emph{Smith v.\ Van Gorkom}, 488
A.2d 858 (Del.\ 1985).
\item \textbf{Aiding and Abetting Breach of Fiduciary Duty} ---
\emph{Malpiede v.\ Townson}, 780 A.2d 1075 (Del.\ 2001).
\item \textbf{Common-Law Fraud} --- Delaware and California (depending
on contracting venue).
\item \textbf{Unjust Enrichment} --- Delaware / California.
\item \textbf{Conversion} --- Delaware / California.
\item \textbf{Tortious Interference with Prospective Economic
Advantage} --- California (Liquidity's principal place of business if
San Francisco-based).
\item \textbf{Civil Conspiracy} --- Delaware / California.
\end{itemize}

\textbf{Forum recommendation:} Federal District Court for the Northern
District of California (likely Liquidity's principal place of business
and the locus of the trade secret misappropriation), with related
Delaware Chancery action for the fiduciary-duty + corporate-governance
claims. Diversity of citizenship and federal-question jurisdiction
both available. Personal jurisdiction over individual Satschel officers
established by their commercial activities directed at Lux (a Delaware
corporation) and / or their California residence.

\textbf{Litigation hold} --- Lux should issue a formal litigation
hold notice to Satschel within 30 days of the SCLA opening offer,
addressed to the corporate entity and named individually to Messrs.\
Choi, Trombley, Church, and Sonsurkar. Failure of Satschel to preserve
records following such notice supports an adverse-inference instruction
at trial.

% ─────────────────────────────────────────────────────────────────
\section*{IV.\ Statement of Facts}

\textit{The following numbered allegations are stated on information
and belief, based on the personal knowledge of the undersigned in his
dual operational role, supplemented by Lux's internal records. Each
will need to be verified, supplemented, or amended through outside-
counsel-supervised diligence prior to any filing.}

\subsection*{A. Background of Lux Industries Inc. and Its IP (2014--2026)}

\textbf{1.} Lux Industries Inc.\ (the predecessor entity Lux Partners
Limited, with continuous corporate identity since 2014) has continuously
developed, since 2014, a blockchain and cryptographic technology stack
including (without limitation) consensus algorithms, post-quantum
cryptographic primitives, threshold and multi-party-computation
protocols, fully-homomorphic-encryption (FHE) precompile architecture,
virtual-machine substrates, EVM derivatives, decentralized-exchange
matching engines, and cross-chain messaging protocols.

\textbf{2.} Lux's continuous direct development investment in this
technology stack, recorded against Lux's books from 2014 to the date
of this memorandum, exceeds \$20 million in direct cost. Fully-loaded
personnel cost over the same period --- senior engineering capacity
not allocated as project cost --- substantially exceeds the direct-
cost figure.

\textbf{3.} Lux's chronological development milestones include:

\begin{itemize}[leftmargin=2em,itemsep=0.1em]
\item \textbf{2014 -- 2017.} Foundational consensus research; early
Quasar consensus prototypes; first cryptographic primitives library;
\texttt{secp256k1}, \texttt{Ed25519}, \texttt{BLS12-381} reference
implementations; HD wallet (BIP-32 / BIP-39) contributions.
\item \textbf{2018 -- 2020.} Active development of the decentralized
exchange (``DEX'') matching engine, including the planet-scale
order-book matching architecture with 1ms finality. Initial EVM fork
from Ava Labs subnet-evm under LGPL inheritance. Threshold-MPC
research (precursors to CGGMP21 and FROST integrations).
\item \textbf{2020 -- 2023.} Quasar consensus stabilization; Pulsar
and Magnetar Apache-2.0-licensed threshold-PQ primitives; FHE
precompile prototypes; precompile registry architecture.
\item \textbf{2024.} Ringtail lattice-signature scheme; post-quantum
standardization push (FIPS 203 / 204 / 205 verifier precompiles
covering ML-KEM, ML-DSA, SLH-DSA); P3Q (Post-Quantum Proof Verifier)
protocol; LEL v1.2 and LRL--PR v1.0 formalized; Warp messaging
protocol stabilization; FHE-based dark-pool matching; multi-chain
architecture (Lux Z-Chain, Q-Chain, X-Chain).
\item \textbf{2025.} LX Suite DEX precompile family (Pool / Oracle /
Router / Hooks at LP-9010 onward), released following Uniswap v4's
hooks model with native on-chain implementation across Lux precompile
addresses; Corona Ring-LWE threshold scheme (successor to Ringtail);
LP-4200 precompile address standardization across the
\texttt{luxfi/precompile/*} registry; multi-arch GPU acceleration under
the closed \texttt{\textasciitilde/work/luxcpp/*} tree (CUDA, Metal,
WebGPU); AI inference precompile (LP-9080 / Compute marketplace).
\item \textbf{2026.} Brand-neutral cleanup across Lux upstream;
relicensing of Pulsar / Magnetar under Apache-2.0 for ecosystem-wide
adoption; ongoing institutional partnership development.
\end{itemize}

\textbf{4.} Lux holds (or has filed for) patents and patent
applications on, without limitation: the Quasar BFT consensus
algorithm; the sampler-bias correction mechanism; the Ringtail and
Corona lattice threshold schemes; the Unified VM Execution Interface
with Pluggable Backends; the Post-Quantum State Verification with
Verkle Proofs protocol; the FHE precompile architecture; the GPU-
accelerated DEX matching engine; and AI-optimized bytecode execution
with JIT compilation.

\subsection*{B. The Lux Licensing Tier Structure}

\textbf{5.} Lux publishes its software under a three-tier licensing
strategy formalized in the internal LICENSING-POLICY (the ``Licensing
Policy''):

\textbf{6.} \textbf{Tier 1 (Commodity / BSD-3-Clause).} Reference
clients, SDKs, common utilities, and other widely-useful but non-
defensible code. Commercial use is free under standard BSD-3
attribution. Examples: \texttt{luxfi/node}, \texttt{luxfi/sdk},
\texttt{luxfi/cli}, \texttt{luxfi/api}, \texttt{luxfi/codec},
\texttt{luxfi/database}, \texttt{luxfi/p2p}, and approximately twenty-
five (25) other modules. Liquidity uses these freely with no Lux
action required.

\textbf{7.} \textbf{Tier 2 (Lux Ecosystem License v1.2, ``LEL'').}
Patent-protected source-available code. Lux publishes source for
inspection, research, audit, and Authorized-Network commercial use,
but reserves commercial rights outside the Lux ecosystem. The license
text at \texttt{\textasciitilde/work/lux/consensus/LICENSE} controls.
Per LEL \S\S 1 and 3, rights are granted only to ``Authorized
Networks'' --- defined to include the Lux Primary Network (NetworkID
1, ChainID 96369), Lux testnets and devnets, and ``L1/L2/L3 chains
descending from the Lux Primary Network.''

\textbf{8.} Per LEL \S 4, forks and competing networks are explicitly
prohibited. Per LEL \S 5, all patent rights are reserved; the
copyright grant does not include any patent license. Per LEL \S 8,
the license terminates automatically on breach. Modules under LEL
include: \texttt{luxfi/consensus}, \texttt{luxfi/crypto},
\texttt{luxfi/mpc}, \texttt{luxfi/threshold}, \texttt{luxfi/ringtail},
\texttt{luxfi/corona}, \texttt{luxfi/fhe}, \texttt{luxfi/precompile}
(and sub-precompiles), \texttt{luxfi/captable}, \texttt{luxfi/broker},
\texttt{luxfi/treasury}, \texttt{luxfi/transfer},
\texttt{luxfi/compliance}, \texttt{luxfi/hsm}, \texttt{luxfi/aml},
\texttt{luxfi/maker}, \texttt{luxfi/warp}, among others.

\textbf{9.} \textbf{Tier 3 (Lux Research License with Patent
Reservation v1.0, ``LRL--PR'').} Research-only commercial license.
Per the LRL--PR text at \texttt{\textasciitilde/work/lux/vm/LICENSE},
Section 2 grants only research-use rights; Section 3 expressly
prohibits any commercial use without separate written license from
Lux; Section 4 expressly reserves all patent claims; Section 5 auto-
licenses validator-node operation only on Lux mainnet and testnets
(no Descending-Chain carve-out); Section 6 contains anti-patent-
filing covenants. Modules under LRL--PR include: \texttt{luxfi/vm},
\texttt{luxfi/staking}, \texttt{luxfi/sampler}, and the Go bindings
of \texttt{luxfi/dex} prior to that repository's move to fully
proprietary status in 2026.

\textbf{10.} \textbf{Tier 4 (Closed Proprietary).} The
\texttt{\textasciitilde/work/luxcpp/*} tree (CUDA / Metal / WebGPU /
FPGA / lattice GPU accelerators, the C++ DEX matching engine, and
related code) is private repository code requiring a separately
negotiated Private-Moat Addendum and commercial license. Liquidity
does not currently use this tier; access in the future requires a
separate written agreement.

\subsection*{C. Satschel's Commercial Use of Lux IP}

\textbf{11.} Satschel (through Satschel, Inc.\ and Liquidity.io LLC)
operates a registered Alternative Trading System (``ATS''),
broker-dealer (``BD''), and transfer-agent (``TA'') business in the
United States. Satschel commercially deploys, in production, on three
live blockchain networks --- Liquid Devnet (ChainID 8675312), Liquid
Testnet (ChainID 8675310), and Liquid Mainnet (ChainID 8675309) ---
the validator binary \texttt{lqd}, built from the
\texttt{liquidityio/node} source tree, which in turn imports and
incorporates Lux modules including \texttt{luxfi/node},
\texttt{luxfi/consensus}, \texttt{luxfi/crypto}, \texttt{luxfi/evm},
\texttt{luxfi/vm}, and dozens of others.

\textbf{12.} Satschel further deploys, in production:

\begin{itemize}[leftmargin=2em,itemsep=0.1em]
\item The Liquidity EVM (\texttt{liquidityio/evm}), which incorporates
\texttt{luxfi/evm} (LGPL-3.0, with \S 6 source-offer obligations),
\texttt{luxfi/vm} (LRL--PR), and \texttt{luxfi/precompile} (LEL with
sub-modules).
\item The Liquidity ATS, BD, and TA backend services, which import
and use \texttt{luxfi/broker}, \texttt{luxfi/treasury},
\texttt{luxfi/transfer}, \texttt{luxfi/captable},
\texttt{luxfi/compliance}, \texttt{luxfi/hsm}, \texttt{luxfi/aml},
\texttt{luxfi/maker}, \texttt{luxfi/crypto}, \texttt{luxfi/mpc},
\texttt{luxfi/threshold}, and \texttt{luxfi/fhe} --- all LEL-licensed.
\item Multiple white-label tenant deployments (including MLC, VCC,
TZERO, ReTokens, and other Satschel-controlled tenants) operating on
the same Lux-licensed substrate.
\item 12,796 SecurityTokens deployed on the Liquid Devnet alone,
each running on top of Lux precompiles for signature verification,
threshold operations, FHE-protected dark-pool matching, and on-chain
identity / compliance.
\end{itemize}

\textbf{13.} The audit underlying this memorandum
(\texttt{LIQUIDITY\_LICENSE\_AUDIT\_2026\_05\_25.md}, co-located in
the Lux Industries Inc.\ legal directory) inventories every Lux module
directly imported by Satschel and assigns each its operative license
tier. Approximately fifteen (15) directly-imported Lux modules carry
LEL or LRL--PR tier licenses requiring an executed commercial license
for Satschel's mode of use.

\textbf{14.} \textbf{No commercial license has ever been executed
between Lux Industries Inc.\ and Satschel.} No side letter, no oral
agreement reduced to writing, no Board resolution authorizing the use
exists in Lux's records to date. Outside counsel should verify the
non-existence of any such instrument through document discovery during
the negotiation or litigation phase.

\textbf{15.} The Liquid blockchain networks operated by Satschel do
not currently activate the Lux Warp messaging precompile (canonical
address \texttt{0x0200000000000000000000000000000000000005}).
Specifically, the \texttt{warpConfig} block is absent from the
genesis JSON of each Liquid EVM chain. Under the LEL Authorized-
Network definition, a chain that does not maintain Lux consensus and
Warp interop does not qualify for the LEL ecosystem-permissive grant.
\textbf{Therefore, even crediting the most generous reading of LEL
\S 1, the Liquid chains do not currently qualify as ``L1/L2/L3 chains
descending from the Lux Primary Network'' as a matter of on-chain
fact.}

\subsection*{D. The Satschel Officer Team and Chain of Knowledge}

\textbf{16.} At all relevant times, Mr.\ Eric Choi served (or now
serves) as Chief Executive Officer of Satschel and Liquidity.io. The
CEO role at any regulated digital-asset venue carries ultimate
operational responsibility for the licensing and intellectual-
property compliance of every technology asset used in production.
Mr.\ Choi's signature appears on the subscription documents for the
\$1.1 billion capital raise; representations made therein regarding
Satschel's technology ownership and licensing posture sound
in his name.

\textbf{17.} At all relevant times, Mr.\ Austin Trombley served as
Chief Strategy Officer of Satschel / Liquidity.io, and prior to that
served as Acting Chief Executive Officer. Mr.\ Trombley is a
blockchain and AI engineer by background; he previously built ATS
infrastructure for Franklin Templeton and served as an eyewitness in
the litigation arising from that engagement. \textbf{Counsel should
review the Franklin Templeton litigation record} --- both for any
pattern-of-conduct evidence relevant here, and for material that may
inform Mr.\ Trombley's likely posture as a witness in any Lux v.\
Satschel proceeding (including any limits on his prior testimony
that may be impeachable).

\textbf{18.} At all relevant times, Mr.\ Coleman Church served as
Chief Executive Officer of Satschel's broker-dealer affiliate. As BD
CEO, Mr.\ Church holds FINRA-regulated responsibility for the BD's
technology and operational integrity. His representations on FINRA
filings and BD compliance attestations are at risk if those filings
do not accurately disclose the Lux licensing status of underlying
BD technology.

\textbf{19.} At all relevant times material to this matter, Mr.\
Jayant Sonsurkar has served as Chief Security Officer of Satschel
and has held maintainer / admin access on the GitHub organization
\texttt{luxfi/*}. The CSO role at any digital-asset venue, and
especially at one carrying ATS/BD/TA registrations, includes (without
limitation):

\begin{itemize}[leftmargin=2em,itemsep=0.1em]
\item Oversight of software-supply-chain integrity, including third-
party-IP review and license compliance.
\item Identification of LEL, LRL--PR, GPL, LGPL, and other
restricted licenses in deployed dependencies.
\item Escalation of licensing mismatches to the CEO, BD CEO, and
General Counsel.
\item Implementation of remediation when a mismatch is identified.
\end{itemize}

\textbf{20.} Mr.\ Sonsurkar's GitHub admin access to \texttt{luxfi/*}
gave him direct, ongoing, contemporaneous visibility into:

\begin{itemize}[leftmargin=2em,itemsep=0.1em]
\item Every Lux \texttt{LICENSE} file as it was authored, amended,
and committed (including the LEL v1.2 December 2025 finalization, the
LRL--PR v1.0 December 2025 finalization, and the relicensing of
specific modules over time).
\item The complete commit history of every Lux module, including
copyright headers, SPDX identifiers, and patent-reservation notices.
\item The identity of every Lux author and the chronological record
of each contribution.
\item Code review and pull-request discussions, including any
discussions of license tier and commercial-use posture.
\end{itemize}

\textbf{21.} On information and belief, Mr.\ Sonsurkar did not, at any
time during the relevant period, raise the LEL / LRL--PR licensing
mismatch internally at Satschel; did not request that Satschel
execute a commercial license with Lux; did not include the
unlicensed-IP status in any FINRA / SEC compliance review; and did
not flag the issue during the \$1.1 billion capital-raise diligence
process. This omission is the central evidentiary fact supporting the
breach-of-fiduciary-duty count (Count IV) and the aiding-and-abetting
count (Count V), and tends to support the Path B factual
characterization described in \S V (intentional misappropriation
rather than mere negligence).

\textbf{22.} On information and belief, Messrs.\ Choi, Trombley, and
Church each had operational and / or strategic visibility into
Satschel's technology stack and the licensing posture of its
underlying components. The factual record will need to develop
through outside-counsel-supervised discovery the precise scope of
each officer's actual knowledge.

\subsection*{E. The Position of Mr.\ Zach Kelling}

\textbf{23.} Mr.\ Zach Kelling is the Founder of both Lux Industries
Inc.\ and Satschel, Inc. He currently serves as Architect and Chief
Cryptographer of Satschel. \textbf{He is not currently the Chief
Executive Officer or Chief Technology Officer of Satschel}, having
previously held those titles in earlier corporate periods (which
this memorandum acknowledges expressly to avoid any inference of
current operational authority on his part).

\textbf{24.} Mr.\ Kelling is a \textbf{minority shareholder} of Lux
Industries Inc., and a \textbf{zero-holder} of any Satschel /
Liquidity.io equity. As such, Mr.\ Kelling does not stand to gain
financially from Satschel's commercial outcome and has not been
positioned to direct Satschel's commercial decisions, licensing
posture, or representations to investors.

\textbf{25.} \textbf{The factual record supports a finding that Mr.\
Kelling has been operationally displaced from Satschel's commercial
direction and has not been a principal actor in the conduct giving
rise to the claims in this memorandum.} The relevant Satschel
officers --- Messrs.\ Choi, Trombley, Church, and Sonsurkar --- have
operated independently of Mr.\ Kelling in the matters at issue. Lux's
litigation strategy ($\S$ XII) should reflect Mr.\ Kelling's distinct
position.

\subsection*{F. Pre-Kelling Technical State and the ACH Exploit (Material Context)}

\textbf{28.} The factual record establishes that, in the period
immediately preceding Mr.\ Zach Kelling's arrival as active Chief
Technology Officer of Satschel, the Satschel ATS platform was
materially non-functional as a regulated alternative trading system.
Specifically, on information and belief:

\begin{itemize}[leftmargin=2em,itemsep=0.1em]
\item The Satschel ATS \textbf{did not implement working time-price
priority order matching} --- the foundational ordering principle of
every regulated US alternative trading system, required by SEC Reg
ATS for any venue that purports to operate as a securities trading
facility.
\item No functioning ATS pipeline existed end-to-end. Critical
components, including order acceptance, matching engine logic,
settlement instructions, and post-trade reporting, were either
non-functional, partially built, or built in a manner that did not
satisfy basic regulatory or commercial requirements.
\item The platform was operationally vulnerable to elementary
security exploits. \textbf{One day prior to Mr.\ Kelling's arrival as
active CTO, the Satschel platform suffered an ACH-funding exploit
that resulted in the theft of approximately \$40,000} via a trivial
attack vector that would have been remediated by even baseline
security review. The exploit underscores both the security posture
of the platform under the pre-Kelling officer team and the absence
of competent CSO oversight (Mr.\ Sonsurkar's role) during that
period.
\end{itemize}

\textbf{29.} Upon his arrival as active CTO, Mr.\ Kelling and his
team (substantially overlapping with the Lux Industries Inc.\
engineering team in personnel and methodology) rebuilt the Satschel
ATS, BD, and TA technology stack from the ground up, importing the
Lux modules described in \S IV(C) as the foundation. \textbf{In
substantial part, the technology that subsequently supported
Satschel's \$1.1 billion capital raise was the Lux IP that Mr.\
Kelling brought across, integrated with the regulatory wrapper.}

\textbf{30.} The technical leadership transition from the pre-
Kelling officer team to the Kelling era is the inflection point that
made Satschel a fundable enterprise. The factual proposition that
Satschel's enterprise value, as represented to investors in the
\$1.1B raise, is materially attributable to Lux-origin technology
contributed by Mr.\ Kelling (and the broader Lux team) is supported
by the contemporaneous technical record.

\textbf{31.} This material context bears on (a) the damages
attribution analysis in \S VI (the share of \$1.1B-valuation
enterprise value attributable to Lux contribution); (b) the
credibility of Mr.\ Trombley as a technology witness in any
subsequent proceeding, given the documented prior technical state;
(c) the breach-of-fiduciary-duty count (\S V, Count IV) and the
absence of CSO oversight evidenced by the ACH exploit; and (d) the
protective treatment of Mr.\ Kelling described in \S XII (his role
was that of remediator and rescuer, not contributor to the
underlying misappropriation).

\subsection*{G. The \$1.1 Billion Capital Raise}

\textbf{26.} In or around the most recent capital raise, Satschel
obtained approximately \$1.1 billion in investor capital at a post-
money valuation in that range. The subscription agreements, private
placement memorandum (``PPM''), and side letters from that raise will
require outside-counsel review for any representations regarding:

\begin{itemize}[leftmargin=2em,itemsep=0.1em]
\item Satschel's ownership of the underlying technology stack.
\item Licensing status of all third-party intellectual property.
\item Related-party transactions including any Lux-Satschel
arrangements.
\item The CSO's compliance attestations.
\item The BD CEO's compliance attestations.
\item Technology disclosures in the risk factors section.
\end{itemize}

\textbf{27.} To the extent the raise documents represent (or omit
material facts regarding) Satschel's ownership or licensed access to
Lux IP, those representations and omissions are independently
actionable under Section 10(b) of the Securities Exchange Act, 15
U.S.C.\ \S 78j(b), and Rule 10b-5, 17 C.F.R.\ \S 240.10b-5. See
Count IX, infra.

% ─────────────────────────────────────────────────────────────────
\section*{V.\ Causes of Action Available to Lux Industries Inc.}

The following causes of action are available based on the factual
record summarized in \S IV. Each is stated with its statutory or
common-law basis, elements, factual support, and applicable damages
framework. Outside counsel should refine the pleading drafts in
light of jurisdiction-specific pleading requirements and the
record developed in pre-suit discovery.

\subsection*{Count I --- Misappropriation of Trade Secrets (Defend
Trade Secrets Act, 18 U.S.C.\ \S 1836)}

\textbf{Elements.} Plaintiff must show: (i) ownership of a trade
secret; (ii) related to a product or service used in interstate
commerce; (iii) that defendant misappropriated; (iv) by improper
means or by acquiring with knowledge of improper acquisition.

\textbf{Factual support.} Lux's unpublished optimizations, integration
know-how, architectural documents, deployment recipes, MPC operational
procedures, sampler-bias remediation methods, FHE operating envelopes,
and pre-publication code in private repositories all qualify as trade
secrets. The Liquid chain bootstrap, the MPC custody integration, and
the FHE dark-pool matching configurations exposed to Satschel through
Mr.\ Sonsurkar's repository access constitute interstate use; Mr.\
Sonsurkar's continuing access without an executed license, coupled
with Satschel's failure to flag the licensing mismatch despite his
duty to do so, supports misappropriation by improper means.

\textbf{Remedies.} Injunctive relief, 18 U.S.C.\ \S 1836(b)(3)(A);
actual damages and unjust enrichment, \S 1836(b)(3)(B); reasonable
royalty in lieu, id.; exemplary damages up to 2$\times$ for willful
and malicious misappropriation, \S 1836(b)(3)(C); attorneys' fees,
\S 1836(b)(3)(D); seizure order if record warrants, \S 1836(b)(2).

\subsection*{Count II --- Copyright Infringement (17 U.S.C.\ \S 501)}

\textbf{Elements.} Plaintiff must show: (i) ownership of valid
copyright; (ii) copying of constituent elements of the work that are
original.

\textbf{Factual support.} Lux owns copyright in every source file
under the \texttt{luxfi/*} GitHub organization. Each Lux author's
contribution vests copyright in Lux under standard assignment
language; copyright notices are present in source headers. Satschel
has copied --- by direct import in \texttt{go.mod}, by container-
image embedding, by binary distribution to validator nodes, and by
on-chain deployment of derivative SecurityTokens --- substantially all
constituent elements of the LEL and LRL--PR-licensed Lux source.

\textbf{Remedies.} Actual damages and infringer's profits, 17 U.S.C.\
\S 504(b); or statutory damages up to \$150,000 per work willfully
infringed, \S 504(c)(2); injunction, \S 502; impoundment and
destruction of infringing copies, \S 503; attorneys' fees, \S 505;
costs, \S 505.

\subsection*{Count III --- Patent Infringement (35 U.S.C.\ \S 271)}

\textbf{Elements.} Plaintiff must show: (i) ownership of valid,
enforceable patent claim; (ii) infringement by making, using,
selling, offering for sale, or importing into the United States
without authority.

\textbf{Factual support.} Lux holds or has filed for patents covering
the inventions enumerated in \S IV(A)(4), supra. Satschel's
operation of validator nodes implementing Quasar consensus, its
deployment of Ringtail / Corona threshold signatures, its use of the
sampler-bias correction mechanism, its operation of the Liquid EVM
substrate implementing the Unified VM Execution Interface, its FHE
precompile activations, and its DEX matching engine each implicate
one or more patent claims. The LRL--PR text \S 4 expressly reserves
the patent rights, so the copyright license that Satschel might
otherwise argue (and which it cannot validly argue) does not include
any patent grant.

\textbf{Remedies.} Reasonable royalty (the statutory floor), 35
U.S.C.\ \S 284; lost profits where established; injunctive relief,
\S 283; treble damages for willful infringement, \S 284; attorneys'
fees in exceptional cases, \S 285; pre- and post-judgment interest.

\subsection*{Count IV --- Breach of Fiduciary Duty (Delaware Law,
against Jayant Sonsurkar individually)}

\textbf{Elements.} As a corporate officer, defendant owes the
corporation duties of care and loyalty. Breach occurs upon failure to
exercise the care of an ordinarily prudent person under similar
circumstances, or by acts that are not in good faith, involve a
conflict of interest, or otherwise constitute gross negligence.

\textbf{Factual support.} Mr.\ Sonsurkar's failure, throughout his
tenure as CSO with admin access to the \texttt{luxfi/*} repositories,
to identify and flag the LEL / LRL--PR licensing mismatch is at
minimum gross negligence in the CSO duty of care. To the extent the
factual record supports knowledge or willful blindness, the conduct
crosses into duty-of-loyalty breach. The duty in this count runs
to Satschel itself (and derivatively to its shareholders, including
investors in the \$1.1B raise); Lux's claim sounds in aiding and
abetting (Count V) rather than direct duty.

\textbf{Remedies.} Damages payable to Satschel (or its creditors /
shareholders in derivative or class capacity); equitable relief
including removal; punitive damages if record warrants.

\subsection*{Count V --- Aiding and Abetting Breach of Fiduciary
Duty (Delaware, against Messrs.\ Choi, Trombley, Church)}

\textbf{Elements.} (i) Existence of a fiduciary relationship; (ii)
breach of fiduciary duty by the fiduciary; (iii) knowing participation
in the breach by the non-fiduciary defendant; (iv) damages.
\emph{Malpiede v.\ Townson}, 780 A.2d 1075 (Del.\ 2001).

\textbf{Factual support.} To the extent Messrs.\ Choi, Trombley, and
Church had actual or constructive knowledge of the licensing mismatch
and continued to direct Satschel's commercial operations notwithstanding,
they aided and abetted Mr.\ Sonsurkar's breach.

\subsection*{Count VI --- Unjust Enrichment}

\textbf{Elements.} (i) Enrichment of the defendant; (ii) at plaintiff's
expense; (iii) absence of justification; (iv) absence of remedy at law.

\textbf{Factual support.} Satschel's commercial deployment of Lux IP
without an executed license enriches Satschel at Lux's expense; the
LEL and LRL--PR expressly negate any justification for commercial use
absent a license. Restitution should measure the value that Satschel
captured, plausibly by reference to (a) the capital raised at \$1.1B
attributable to the unlicensed IP, (b) the gross-revenue royalty Lux
would have charged at arm's-length, or (c) the avoided build-cost
Satschel saved by relying on Lux's twelve-year investment.

\subsection*{Count VII --- Conversion}

\textbf{Elements.} (i) Plaintiff's interest in identifiable property;
(ii) defendant's intentional interference inconsistent with plaintiff's
rights.

\textbf{Factual support.} Satschel's deployment of Lux IP without
authorization, particularly the patent-protected modules under LRL--PR,
constitutes intentional interference with Lux's exclusive rights.

\subsection*{Count VIII --- Tortious Interference with Prospective
Economic Advantage}

\textbf{Elements.} (i) Existence of a prospective economic relationship
with third parties having probability of future economic benefit to
plaintiff; (ii) defendant's knowledge of the relationship; (iii)
intentional and wrongful interference; (iv) actual disruption; (v)
damages.

\textbf{Factual support.} Lux's commercial-license pipeline with
third-party ecosystem partners is impaired by Satschel's unlicensed
use of the same modules; Satschel's offering of effectively
``competitor'' services on Lux IP without paying Lux's customary
license fee depresses the market price third parties will pay for
their own licenses. Satschel's white-label tenants (MLC, VCC, TZERO,
ReTokens) constitute particular instances of this interference.

\subsection*{Count IX --- Securities Fraud (Rule 10b-5)}

\textbf{Elements.} (i) Material misrepresentation or omission; (ii)
scienter; (iii) connection with the purchase or sale of a security;
(iv) reliance; (v) economic loss; (vi) loss causation.

\textbf{Factual support.} To the extent the \$1.1B raise was supported
by representations regarding Satschel's IP ownership or licensed
access that omit or mischaracterize the unlicensed status of the Lux
technology, the elements are met. Satschel officers had scienter
through Mr.\ Sonsurkar's repository access; investors relied on
the representations; the omission was material to a reasonable investor
in a technology-backed venture; and loss causation is established by
the value attributable to the misrepresented or omitted IP rights.
Lux's standing depends on whether Lux itself purchased Satschel
securities (likely not); however, Lux may have collateral standing as
a third-party beneficiary and may also pursue regulator-referral
remedies. \textbf{Counsel should evaluate this count carefully ---
its primary utility is as collateral leverage rather than as a
direct claim by Lux.}

\subsection*{Count X --- Lanham Act False Advertising (15 U.S.C.\
\S 1125(a))}

\textbf{Elements.} (i) False or misleading statement about defendant's
goods or services; (ii) materiality; (iii) actual deception or
tendency to deceive; (iv) commerce; (v) injury.

\textbf{Factual support.} Satschel's marketing materials describe its
technology stack (post-quantum cryptography, FHE-based dark-pool
matching, MPC custody, multi-chain architecture, Quasar-derived
consensus) in terms that ascribe ownership or proprietary control to
Satschel without disclosing the Lux origin or the conditional /
unlicensed nature of the access. To the extent reasonable consumers
are deceived as to origin or quality (or as to Satschel's authority
to provide the services), Lanham Act false-advertising claims attach.

\subsection*{Count XI --- Common-Law Fraud}

\textbf{Elements.} (i) False representation; (ii) knowledge of falsity
or reckless disregard; (iii) intent to induce reliance; (iv)
justifiable reliance; (v) injury.

\textbf{Factual support.} As to Lux directly: to the extent Satschel
officers represented to Lux personnel (Mr.\ Kelling, Mr.\ Bin, or
others) that licensing arrangements would be formalized in due course
and Lux relied to its detriment by continuing to contribute engineering
work without an executed agreement, common-law fraud may lie. The
factual development here will depend on specific representations made
to specific Lux personnel; this is a count that needs careful pleading
support.

\subsection*{Count XII --- Civil Conspiracy}

\textbf{Elements.} (i) Two or more persons; (ii) an unlawful act;
(iii) by unlawful means; (iv) an agreement; (v) damages.

\textbf{Factual support.} To the extent the Satschel officer team
(Messrs.\ Choi, Trombley, Church, and Sonsurkar) acted in concert
in the commercial deployment of unlicensed Lux IP and the raising of
capital on the back of that deployment, civil conspiracy lies. This
count provides joint and several liability and discovery latitude.

% ─────────────────────────────────────────────────────────────────
\section*{VI.\ Damages Analysis}

The damages owed to Lux Industries Inc.\ scale across three scenarios.
Counsel should treat these as a framework; ranges will tighten with
discovery.

\subsection*{Small Scenario (Path A --- Negligent Posture, Settlement)}

The factual record supports a Path A characterization (CSO duty-of-care
failure / negligent oversight; no proven intent). Damages are the
back-license amount plus a reasonable forward-looking grant:

\begin{center}
\begin{tabular}{lr}
\toprule
\textbf{Item} & \textbf{Amount} \\
\midrule
Back-license value (12 yr at \$2.5M/yr base) & \$30M \\
Pre-judgment interest (5\%, simple) & \$10M \\
Royalty component (1.5\% of estimated GR) & \$10M \\
Forward-looking license value (NPV) & \$30M \\
\midrule
\textbf{Subtotal --- Small Scenario} & \textbf{\$80M} \\
\bottomrule
\end{tabular}
\end{center}

\subsection*{Medium Scenario (Path B --- Mixed-Knowledge, Contested)}

The factual record develops to show that one or more officers had
actual knowledge of the licensing mismatch and the conduct continued
unabated. Statutory multipliers apply selectively. Damages
approximately 2.5--3$\times$ the Small Scenario:

\begin{center}
\begin{tabular}{lr}
\toprule
\textbf{Item} & \textbf{Amount} \\
\midrule
Small-Scenario base & \$80M \\
DTSA 2$\times$ exemplary on willful misappropriation & +\$30M \\
Patent willfulness component (selective trebling) & +\$40M \\
Disgorgement on \$1.1B-raise attribution (5\%) & +\$55M \\
Attorneys' fees and costs & +\$10M \\
\midrule
\textbf{Subtotal --- Medium Scenario} & \textbf{\$215M} \\
\bottomrule
\end{tabular}
\end{center}

\subsection*{Big Scenario (Path B Fully Litigated, Adverse Officer Record)}

Factual record establishes intentional misappropriation by multiple
officers; full statutory multipliers apply; capital-raise attribution
expanded; Lanham Act and securities-fraud counts contribute:

\begin{center}
\begin{tabular}{lr}
\toprule
\textbf{Item} & \textbf{Amount} \\
\midrule
Base + interest + royalty & \$80M \\
DTSA full 2$\times$ exemplary & +\$80M \\
Patent full 3$\times$ trebling on willfulness & +\$120M \\
Disgorgement on \$1.1B-raise attribution (15--20\%) & +\$200M \\
Lanham Act + Lanham fees & +\$10M \\
Securities fraud component (if Lux can establish standing) & +\$50M \\
Attorneys' fees, costs, pre-judgment interest & +\$50M \\
\midrule
\textbf{Subtotal --- Big Scenario} & \textbf{\$590M} \\
\bottomrule
\end{tabular}
\end{center}

\subsection*{Opportunity-Cost Damages: 90-Day Bandwidth Diversion}

\textbf{The damages framework above does not yet capture a distinct
and significant category of harm: the opportunity cost to Lux
Industries Inc.\ of having Mr.\ Kelling's --- and the broader Lux
engineering team's --- bandwidth diverted to Satschel's platform
during a critical 90-day window. The data is unambiguous.}

\textbf{Quantification of contribution to Satschel (rolling 90-day
window ending on the memorandum date).} Audit of the
\texttt{liquidityio/*} GitHub organization for commits attributable
to Mr.\ Kelling (across the verified author identities
\texttt{zatsch}, \texttt{z@liquidity.io}, \texttt{z@satschel.com},
and Kelling-signed) yields the following:

\begin{center}
\begin{small}
\begin{tabular}{lrrr}
\toprule
\textbf{Repo} & \textbf{Zach commits} & \textbf{Total commits} & \textbf{Zach \%} \\
\midrule
\texttt{liquidityio/ats}        & 806    & 1{,}426 & 57\% \\
\texttt{liquidityio/bd}         & 388    & 751     & 52\% \\
\texttt{liquidityio/ta}         & 208    & 364     & 57\% \\
\texttt{liquidityio/mpc}        & 31     & 95      & 33\% \\
\texttt{liquidityio/evm}        & 23     & 32      & 72\% \\
\texttt{liquidityio/fhe}        & 10     & 14      & 71\% \\
\texttt{liquidityio/iam}        & 27     & 29      & 93\% \\
\texttt{liquidityio/node}       & 338    & 380     & 89\% \\
\texttt{liquidityio/operator}   & 354    & 406     & 87\% \\
\texttt{liquidityio/universe}   & 1{,}436 & 1{,}920 & 75\% \\
\texttt{liquidityio/contracts}  & 72     & 176     & 41\% \\
\texttt{liquidityio/dex}        & 24     & 32      & 75\% \\
\texttt{liquidityio/platform}   & 36     & 51      & 71\% \\
\texttt{liquidityio/superadmin} & 243    & 338     & 72\% \\
\texttt{liquidityio/exchange}   & 522    & 1{,}170 & 45\% \\
\texttt{liquidityio/console}    & 28     & 38      & 74\% \\
\texttt{liquidityio/ui}         & 27     & 143     & 19\% \\
\texttt{liquidityio/sdk}        & 30     & 30      & 100\% \\
\texttt{liquidityio/sdk-go}     & 9      & 9       & 100\% \\
\texttt{liquidityio/cli}        & 114    & 117     & 97\% \\
\midrule
\textbf{Total (90-day window)} & \textbf{4{,}726} & \textbf{7{,}521} & \textbf{63\%} \\
\bottomrule
\end{tabular}
\end{small}
\end{center}

\textbf{Code-churn magnitude across the largest repositories:}

\begin{center}
\begin{small}
\begin{tabular}{lrr}
\toprule
\textbf{Repo} & \textbf{Files touched by Zach} & \textbf{Lines (added + deleted)} \\
\midrule
\texttt{liquidityio/ats}        & 64{,}899 & $\sim$21.4M \\
\texttt{liquidityio/bd}         & 55{,}878 & $\sim$5.9M \\
\texttt{liquidityio/ta}         & 69{,}321 & $\sim$35.5M \\
\texttt{liquidityio/universe}   & 11{,}493 & $\sim$760K \\
\texttt{liquidityio/exchange}   & 41{,}259 & $\sim$3.9M \\
\texttt{liquidityio/superadmin} & 24{,}688 & $\sim$3.2M \\
\texttt{liquidityio/node}       & 26{,}755 & $\sim$7.4M \\
\texttt{liquidityio/operator}   & 975      & $\sim$71K \\
\bottomrule
\end{tabular}
\end{small}
\end{center}

(Note: file counts include vendored-dependency churn; the unique-
authored-line figure is materially lower than the raw line-count but
of the same order of magnitude. The point is the magnitude, which
is unambiguous --- Mr.\ Kelling committed code on \textbf{55 distinct
days in the ATS repository alone} during the 90-day window, i.e.,
six days a week of continuous output.)

\textbf{Quantification of contribution to Lux upstream in the same
90-day window.} Audit of the \texttt{luxfi/*} GitHub organization for
commits attributable to Mr.\ Kelling in the same period yields:

\begin{center}
\begin{small}
\begin{tabular}{lrr}
\toprule
\textbf{Lux repo} & \textbf{Zach commits} & \textbf{Total commits (Lux team)} \\
\midrule
\texttt{luxfi/node}             & \textbf{0} & 179 \\
\texttt{luxfi/consensus}        & \textbf{0} & 104 \\
\texttt{luxfi/crypto}           & \textbf{0} & 64 \\
\texttt{luxfi/evm}              & \textbf{0} & 99 \\
\texttt{luxfi/vm}               & \textbf{0} & 27 \\
\texttt{luxfi/precompile}       & \textbf{0} & 98 \\
\midrule
\textbf{Total (90-day window)}  & \textbf{0} & 571 \\
\bottomrule
\end{tabular}
\end{small}
\end{center}

\textbf{The pattern is unambiguous.} During the 90-day window in
which Mr.\ Kelling committed 4,726 lines of work to Satschel ---
contributing the substantial majority of all engineering output in
the company's largest repositories --- he committed \textbf{zero}
to the upstream Lux repositories where he serves as Founder, Chief
Cryptographer, and minority shareholder. The Lux upstream during
that period was maintained by other Lux engineers without the
benefit of his attention.

\textbf{Value-creation rate translation.} Satschel's reported
post-money valuation following its recent capital raise is
approximately \$1.1 billion. To the extent the platform that
supported that raise was substantially built or rebuilt by Mr.\
Kelling during the 90-day window analyzed here (and the
contemporaneous git record establishes that it was), the implied
value-creation rate of his bandwidth is approximately
\textbf{\$1.1 billion / 90 days = \$12.2 million per day} of
diverted attention. \textbf{The opportunity cost to Lux Industries
Inc.}\ ---measured by what Lux's own roadmap would have produced
during the same window had Mr.\ Kelling's attention been available
to Lux upstream --- \textbf{is plausibly bounded above by the same
rate, and is bounded below by the proportional share of Lux's
existing enterprise trajectory attributable to founder-level
engineering input}.

\textbf{Damages calculation.} A reasonable measure of opportunity-
cost damages is:

\begin{center}
\begin{tabular}{lr}
\toprule
\textbf{Component} & \textbf{Range} \\
\midrule
Bandwidth diversion at 50\% of value-creation rate (90 days) & \$500M -- \$600M \\
Bandwidth diversion at 25\% rate & \$250M -- \$300M \\
Bandwidth diversion at 10\% rate (conservative) & \$100M -- \$120M \\
\midrule
\textbf{Conservative opportunity-cost claim} & \textbf{\$100--300M} \\
\bottomrule
\end{tabular}
\end{center}

This opportunity-cost component is independent of, and additive to,
the IP-misappropriation damages set out in the Small / Medium / Big
scenarios above. The claim sounds in tortious-interference-with-
contract (Lux's implicit / fiduciary expectation of Mr.\ Kelling's
attention as Founder), unjust enrichment (Satschel captured value
that should have flowed to Lux), and breach of fiduciary duty by
Satschel officers (who knowingly directed Mr.\ Kelling's bandwidth
to Satschel deliverables in lieu of Lux deliverables despite his
dual obligations).

\textbf{This is the single most striking quantitative fact in the
memorandum} and should be lead-in evidence for Mr.\ Dupont in any
opening negotiation. Satschel built itself to a \$1.1B valuation in
the 90 days Mr.\ Kelling's bandwidth was directed there; Lux was
deprived of that same bandwidth at exactly the rate Satschel
captured it.

\subsection*{Damages Summary (Including Opportunity Cost)}

\begin{center}
\begin{tabular}{lrr}
\toprule
\textbf{Scenario} & \textbf{IP damages} & \textbf{+ Opp.\ Cost} \\
\midrule
Small (Path A, settlement context) & \$80M & \$180--380M \\
Medium (Path B, mixed-record litigation) & \$215M & \$315--515M \\
Big (Path B, full-record litigation w/ multipliers) & \$590M & \$690--890M \\
\bottomrule
\end{tabular}
\end{center}

The Big-Scenario range is supported by the Georgia-Pacific royalty
framework (1.5--3.0\% of gross trading + settlement revenue applied
to twelve years of unlicensed use, with multipliers), the patent
willfulness ceiling (3$\times$ under \S 284), the DTSA exemplary
multiplier (2$\times$ under \S 1836(b)(3)(C)), and the unjust-
enrichment recoupment against the recent capital raise.

% ─────────────────────────────────────────────────────────────────
\section*{VII.\ Trading-Fee Revenue Projections (User-Scale Analysis)}

To assist Mr.\ Dupont's deal modeling, the following projects
Liquidity's enterprise-value trajectory at four user-base milestones
and computes Lux's economic interest at each, under the proposed
10\% equity stake (plus optional royalty under alternative
structures). Three Average Revenue Per User (``ARPU'') tiers are
modeled --- conservative (\$50/yr), base (\$100/yr), and premium
(\$200/yr) --- to bracket the realistic range for a tokenized-asset
exchange + ATS/BD/TA combined offering.

\subsection*{Annual Revenue, Implied Enterprise Value, and Lux 10\%
Stake Value}

\begin{center}
\begin{small}
\begin{tabular}{lrrrr}
\toprule
\textbf{Scenario} & \textbf{Users} & \textbf{Revenue} & \textbf{Enterprise Value} & \textbf{Lux 10\% Stake} \\
\midrule
\multicolumn{5}{l}{\textit{Conservative ARPU (\$50/year)}} \\
Stage 1 & 100K & \$5M & \$50M (10$\times$) & \$5M \\
Stage 2 & 1M & \$50M & \$500M (10$\times$) & \$50M \\
Stage 3 & 10M & \$500M & \$5B (10$\times$) & \$500M \\
Stage 4 & 100M & \$5B & \$50B (10$\times$) & \$5B \\
\midrule
\multicolumn{5}{l}{\textit{Base ARPU (\$100/year)}} \\
Stage 1 & 100K & \$10M & \$150M (15$\times$) & \$15M \\
Stage 2 & 1M & \$100M & \$1.5B (15$\times$) & \$150M \\
Stage 3 & 10M & \$1B & \$15B (15$\times$) & \$1.5B \\
Stage 4 & 100M & \$10B & \$150B (15$\times$) & \textbf{\$15B} \\
\midrule
\multicolumn{5}{l}{\textit{Premium ARPU (\$200/year)}} \\
Stage 1 & 100K & \$20M & \$400M (20$\times$) & \$40M \\
Stage 2 & 1M & \$200M & \$4B (20$\times$) & \$400M \\
Stage 3 & 10M & \$2B & \$40B (20$\times$) & \$4B \\
Stage 4 & 100M & \$20B & \$400B (20$\times$) & \textbf{\$40B} \\
\bottomrule
\end{tabular}
\end{small}
\end{center}

\textbf{Lux's 10\% equity, retained through the growth curve, is
worth between \$5 million (Stage 1, Conservative ARPU) and \$40
billion (Stage 4, Premium ARPU).} The mid-case (Base ARPU at 100M
users) implies \$15 billion of Lux equity value --- which is the
order of magnitude at which the user-base projection meaningfully
exceeds any plausible damages award in litigation.

\subsection*{Trading-Fee Royalty Alternative (Approach C scenario)}

Under a hybrid structure (5\% equity + 5 basis-points-of-volume
royalty), Lux's royalty stream at each user-base scale:

\begin{center}
\begin{small}
\begin{tabular}{lrrr}
\toprule
\textbf{Scenario} & \textbf{Implied Volume} & \textbf{5 bps Royalty} & \textbf{Capped at 10\% of Rev} \\
\midrule
Stage 1 (100K users) & \$5--20B & \$2.5--10M & \$0.5--2M \\
Stage 2 (1M users)   & \$50--200B & \$25--100M & \$5--20M \\
Stage 3 (10M users)  & \$500B--2T & \$250M--1B & \$50--200M \\
Stage 4 (100M users) & \$5T--20T  & \$2.5B--10B & \$0.5--2B \\
\bottomrule
\end{tabular}
\end{small}
\end{center}

The royalty stream provides current-period cash that the pure equity
does not; at Stage 3+ it materially exceeds the equity-only NPV at
similar growth assumptions.

\subsection*{Win-Sharing at Larger Equity Stakes}

Aspirational scenarios (a 50/50 Lux-Satschel partnership, or a
Lux-led majority recapitalization following a successful Big-Scenario
litigation) would scale the figures above proportionately --- Lux's
50\% stake at Stage 4 Premium would be worth on the order of \$200
billion. These scenarios should be understood as ceiling cases for
Mr.\ Dupont's negotiating posture, not as expected outcomes.

% ─────────────────────────────────────────────────────────────────
\section*{VIII.\ Strategic Partnership Upside}

Lux brings substantial commercial value to a partnership-resolution
of this dispute --- value that is independent of the IP settlement
itself, and that should be quantified and surfaced in any
negotiation as the upside Mr.\ Dupont can offer in exchange for
favorable settlement terms.

\subsection*{Real-World Asset Pipeline}

Lux operates the infrastructure layer behind multi-billion-dollar
RWA tokenization pipelines across multiple sovereign chains. Post-
SCLA, Liquidity is the natural US-regulated listing destination for
these flows:

\begin{itemize}[leftmargin=1.5em,itemsep=0.1em]
\item Tokenized US securities + private equity (Carta ingestion);
\item Tokenized real estate / NNN-SPV-wrapped property exposures
(the \texttt{net\_lease\_spe} asset class);
\item Tokenized commodities (gold, silver, FX, energy);
\item Tokenized fund LP interests (PE / VC, hedge funds);
\item Tokenized fixed-income (US Treasuries, corporate debt,
sovereign debt under Reg D / Reg S);
\item Tokenized music + IP royalty streams;
\item Tokenized art and collectibles (\$1M+ lots).
\end{itemize}

\textbf{Conservative estimate of incremental Liquidity deal flow
under the SCLA: \$5--15 billion in nominal asset value over the first
12--24 months post-close.}

\subsection*{Banking Licenses and Financial Infrastructure (via SFPB
Partnership)}

Lux operates a comprehensive licensed banking and financial
infrastructure stack via its SF Private Bank (``SFPB'') partnership.
The combined regulatory footprint covers multiple jurisdictions:

\begin{itemize}[leftmargin=1.5em,itemsep=0.1em]
\item \textbf{United States} --- federal banking partner; FDIC-
insured rails.
\item \textbf{Canada} --- chartered banking partner.
\item \textbf{Singapore} --- MAS-licensed (Capital Markets Services
+ banking) partner.
\item \textbf{New Zealand} --- RBNZ-licensed deposit-taker partner.
\item \textbf{Sweden} --- Finansinspektionen-supervised; EU passport.
\item Additional jurisdictions (UK FCA, UAE VARA, Switzerland FINMA,
Australia AUSTRAC) under partnership coverage.
\end{itemize}

Services Lux can deliver to Liquidity under the SCLA include
Banking-as-a-Service (BaaS), Payment Service Provider (PSP) rails,
and \textbf{a custody solution that materially undercuts the
Fireblocks and BitGo incumbents on price} (built on the same MPC
infrastructure already in Liquidity production).

\textbf{Estimated independent cost for Satschel to acquire equivalent
global license suite: \$57--140M in direct regulatory and capital
expenditure, plus 5--7 years of timeline.} This regulatory access
alone exceeds the floor settlement value of any negotiating tier.

\subsection*{International Distribution: M7 and 1B+ End-User Reach}

Through Hanzo Inc., the broader corporate family, and the Lux
ecosystem partnerships, the M7 global initiatives provide
distribution access to more than 1 billion end-users across consumer,
institutional, and retail channels. Post-SCLA, Lux's international
counterparty network operates as a Reg-S-compatible distribution
channel for Liquidity-listed assets, with cross-network Warp messaging
providing the on-chain attestation rail. This expands Liquidity's
investor base by a meaningful multiple beyond domestic US channels.

\subsection*{Hanzo Inc.\ Portfolio Precedents}

Hanzo Inc., the venture studio under which the Lux Industries
relationship has historically operated, has applied the IP-in-kind
founder-equity model successfully to multiple spawned companies
that have since reached unicorn or near-unicorn valuations:

\begin{itemize}[leftmargin=1.5em,itemsep=0.1em]
\item \textbf{Damon Motorcycles} --- electric / smart-motorcycle company.
\item \textbf{Triller} --- short-form video / social platform.
\item \textbf{Bellabeat} --- women's health wearables.
\item \textbf{Unicorn Gold} --- precious-metals tokenization /
investment platform.
\item \textbf{Lux Network (Lux Industries Inc.)} --- itself a Hanzo-
spawn outcome.
\item \textbf{Zoo Labs Foundation / Zoo Network} --- decentralized AI
research network.
\end{itemize}

\textbf{The model proposed for Satschel is not novel}; it has produced
multiple unicorn outcomes in the same corporate family. The retained-
parent-equity range across the portfolio runs 10--40\%, and the 10\%
ask in Satschel sits at the bottom of that internal range.

\subsection*{Ongoing IP Pipeline}

Beyond the modules in production, Lux is actively developing further
IP that becomes available to Satschel under any executed SCLA:
AI inference + compute precompiles; advanced ZK rollup substrates;
post-quantum hybrid schemes; threshold MPC advances (dynamic
re-sharing); confidential-execution attestation (TEE / SEV-SNP / TDX
binding to consensus). \textbf{Under a perpetual ecosystem license,
Satschel receives every future Lux IP contribution as it ships,
without renegotiation or additional fees.}

\subsection*{The \$100B--\$1T Target Institution}

The combined Lux + Liquidity vision is the construction of a
\$100-billion-to-\$1-trillion-class institution at the intersection
of tokenized real-world assets, post-quantum-secure infrastructure,
regulated US securities markets, and global distribution. The Lux
deal team should treat the negotiation not as a zero-sum settlement
of a past dispute, but as the founding-document of a partnership
that creates value materially in excess of the negotiating range.

% ─────────────────────────────────────────────────────────────────
\section*{IX.\ Wait vs.\ Settle Analysis}

\textbf{The question Mr.\ Dupont asked is: ``Is it better to wait?
Liquidity's equity becomes more valuable, and damages we are owed
grow.''} The legal analysis is more nuanced than the simple
formulation suggests.

\subsection*{Forces Favoring Delay}

\begin{itemize}[leftmargin=1.5em,itemsep=0.1em]
\item \textbf{Equity appreciation}: Liquidity's enterprise value
likely grows; Lux's 10\% appreciates with it.
\item \textbf{Royalty accrual}: each year of continued unlicensed
use adds approximately \$5--10M to the Small-Scenario base.
\item \textbf{Capital-raise disgorgement}: future raises expand the
disgorgement universe (the \$1.1B is the floor; a \$5B Series B
multiplies the relevant attribution number).
\item \textbf{Pattern-of-conduct evidence}: longer the conduct
continues, the stronger the willfulness narrative.
\end{itemize}

\subsection*{Forces Favoring Settlement Now}

\begin{itemize}[leftmargin=1.5em,itemsep=0.1em]
\item \textbf{Statutory caps}: DTSA multiplier caps at 2$\times$;
patent multiplier caps at 3$\times$. These do not grow with delay.
\item \textbf{Dilution risk}: every Satschel preferred round dilutes
Lux's relative ownership absent strong anti-dilution protection.
\item \textbf{Operational risk}: Satschel could engineer around the
LEL modules over time, reducing both damages and the future-value
narrative.
\item \textbf{Evidence preservation}: longer the delay, greater the
risk of document destruction (countered by litigation hold).
\item \textbf{Tolling vs.\ DTSA limitations}: 3-year statute of
limitations from discovery; counsel should consider whether to file a
protective complaint or obtain a tolling agreement.
\item \textbf{Lux's operating runway}: a \$10M+ cash settlement
component funds Lux operations now; deferred value does not.
\item \textbf{Negotiating leverage compounds in our favor short-term}
but evaporates if Liquidity hits trouble (regulatory action, market
downturn, or executive scandal at Satschel).
\item \textbf{Public visibility}: Satschel's pre-IPO / Series B /
exam-cycle timing makes the next 6--12 months the highest-leverage
window for Lux.
\end{itemize}

\subsection*{Recommendation}

\textbf{Settle now, with strong protective terms.} The optimal
structure is:

\begin{enumerate}[leftmargin=1.5em,itemsep=0.2em]
\item Significant cash component (\$10M floor; higher preferable)
into Lux primary investment;
\item 10\% Satschel common (or higher) with weighted-average broad-
based anti-dilution;
\item Pro-rata participation right on all future preferred rounds
(maintains Lux's percentage through dilution);
\item Tag-along and co-sale rights (Lux participates in any
liquidity event);
\item Mandatory 1+1 board composition (1 Lux director + 1
independent director added to the existing 3-person board);
\item 99-year exclusive USA-market license to Liquidity, conditional
on Satschel maintaining all SCLA terms;
\item Patent grant on all current and future Lux patents covering
the licensed IP.
\end{enumerate}

\textbf{This captures the upside of growth} (through equity
appreciation under anti-dilution) \textbf{without the risk of
delayed enforcement.} The statutory damages multipliers do not grow
with delay; the equity does. Locking in equity now with strong anti-
dilution captures the future-growth narrative without the operational
risks of continued reliance on a deteriorating informal arrangement.

\textbf{Litigation should be held in reserve, not pursued} ---
litigation extracts a one-time damages payment but forecloses the
ongoing partnership upside. A clean settlement that includes equity
participation is strictly better than a litigated damages award,
\emph{provided} the settlement terms are themselves favorable.

\textbf{Counsel should not deviate from this analysis} unless the
diligence record develops new evidence that materially alters either
the damages range (Big-Scenario evidence emerges) or the partnership-
upside range (Liquidity's trajectory changes materially).

% ─────────────────────────────────────────────────────────────────
\section*{X.\ Settlement Triangulation (Cash / Equity / Features Matrix)}

Rather than propose a single settlement, this section presents the
triangulation matrix the Lux deal team can use to negotiate. The
floor and ceiling reflect the damages analysis (\S VI) and the
partnership upside (\S VIII). Mr.\ Dupont is in the lead on selecting
the operating tier.

\subsection*{Settlement Sizing Tiers}

\begin{center}
\begin{small}
\begin{tabular}{lrrrr}
\toprule
\textbf{Tier} & \textbf{Cash} & \textbf{Equity} & \textbf{Lux Nominal Value} & \textbf{Implied Posture} \\
\midrule
Floor (S--)       & \$5M   & 5\%   & \$60M   & Soft settle; preserve relationship \\
Small (S)         & \$10M  & 10\%  & \$120M  & Approach A baseline (recommended) \\
Medium (M)        & \$20M  & 15\%  & \$185M  & Path B leverage applied softly \\
Big (B)           & \$30M  & 25\%  & \$305M  & Full Path B post-investigation \\
Aspirational (A+) & \$50M  & 35\%  & \$435M  & Litigation-credible ceiling \\
50/50 (max)       & \$100M & 50\%  & \$650M  & Full partnership re-cap \\
\bottomrule
\end{tabular}
\end{small}
\end{center}

\subsection*{Cash vs.\ Equity Mix Sensitivity}

At any given nominal value, the cash/equity mix is negotiable. Cash
provides current-period Lux operating runway; equity provides upside
participation. The following sensitivity table fixes Lux's settlement
value at the Small-Tier \$120M floor and shows the negotiating mix
options:

\begin{center}
\begin{tabular}{lll}
\toprule
\textbf{Mix} & \textbf{Cash} & \textbf{Equity at \$1.1B} \\
\midrule
All-cash & \$120M & 0\% \\
Heavy cash & \$60M & 5.5\% \\
Balanced (recommended) & \$10M & 10\% \\
Light cash & \$5M & 10.5\% \\
All-equity & \$0 & 11\% \\
\bottomrule
\end{tabular}
\end{center}

The \textbf{balanced mix at \$10M cash + 10\% equity is recommended}
on two grounds: (a) the \$10M cash provides material Lux operating
runway without overburdening Satschel's post-raise balance sheet,
and (b) the 10\% equity captures the partnership upside identified
in \S VII. Higher cash components are achievable if Satschel
prefers to deploy raised capital as settlement rather than retain it.

\subsection*{Feature Unlock by Settlement Tier}

The level of forward-looking value Lux is willing to bring as
aligned-partner scales with the settlement tier. At the Floor (S--)
tier, Lux contributes only the historical IP license. At the Big
(B) tier or above, Lux's full ecosystem and infrastructure are
unlocked:

\begin{center}
\begin{small}
\begin{tabular}{p{4.5cm} ccccc}
\toprule
\textbf{Forward-value feature} & \textbf{S--} & \textbf{S} & \textbf{M} & \textbf{B} & \textbf{A+} \\
\midrule
LEL/LRL--PR perpetual license & $\checkmark$ & $\checkmark$ & $\checkmark$ & $\checkmark$ & $\checkmark$ \\
Patent grant on covered IP & $\checkmark$ & $\checkmark$ & $\checkmark$ & $\checkmark$ & $\checkmark$ \\
Maintenance \& support of upstream OSS & & $\checkmark$ & $\checkmark$ & $\checkmark$ & $\checkmark$ \\
99-year exclusive USA market license & & $\checkmark$ & $\checkmark$ & $\checkmark$ & $\checkmark$ \\
1 Lux director + 1 independent seat & & $\checkmark$ & $\checkmark$ & $\checkmark$ & $\checkmark$ \\
RWA pipeline routing into Liquidity & & $\checkmark$ & $\checkmark$ & $\checkmark$ & $\checkmark$ \\
SFPB banking-license access (BaaS, PSP) & & & $\checkmark$ & $\checkmark$ & $\checkmark$ \\
Custody (Fireblocks-undercutting) & & & $\checkmark$ & $\checkmark$ & $\checkmark$ \\
International M7 distribution channels & & & & $\checkmark$ & $\checkmark$ \\
Closed-source (\texttt{luxcpp}) GPU access & & & & & $\checkmark$ \\
Co-marketing + brand alignment & & & & $\checkmark$ & $\checkmark$ \\
Joint go-to-market programs & & & & $\checkmark$ & $\checkmark$ \\
\bottomrule
\end{tabular}
\end{small}
\end{center}

\textbf{Mr.\ Dupont can use this matrix in negotiation directly}:
``Lux is willing to bring [feature set] at [tier], or alternatively
to bring [smaller feature set] at [smaller tier].'' Satschel chooses
its tier; Lux's deliverables scale accordingly.

\subsection*{Damages-vs-Settlement Triangulation}

\begin{center}
\begin{tabular}{lrr}
\toprule
\textbf{Outcome path} & \textbf{Lux receives} & \textbf{Satschel pays} \\
\midrule
Settlement at Floor (S--) & \$60M nominal & \$5M cash + 5\% equity \\
Settlement at Small (S, recommended) & \$120M nominal & \$10M cash + 10\% equity \\
Settlement at Medium (M) & \$185M nominal & \$20M cash + 15\% equity \\
Settlement at Big (B) & \$305M nominal & \$30M cash + 25\% equity \\
Settlement at Aspirational (A+) & \$435M nominal & \$50M cash + 35\% equity \\
Litigation: Small-Scenario damages (Path A) & \$80M damages & \$80M cash, no equity \\
Litigation: Medium-Scenario damages (Path B mixed) & \$215M damages & \$215M, no equity \\
Litigation: Big-Scenario damages (Path B full) & \$590M damages & \$590M, no equity \\
\bottomrule
\end{tabular}
\end{center}

\textbf{Key insight}: a settled outcome at the Small tier (\$120M
nominal) is materially less expensive to Satschel than the Small-
Scenario litigated damages (\$80M cash), \textit{provided the equity
component is properly valued}. Litigation extracts cash without
preserving the partnership relationship and without giving Satschel
the forward-looking benefits in \S VIII. \textbf{Settlement
dominates litigation at every tier for both sides.}

% ─────────────────────────────────────────────────────────────────
\section*{XI.\ Mandatory Terms (Lux's Floor in Any Settlement)}

The following terms are non-negotiable. Any settlement at any tier
must include each of them. Mr.\ Dupont and Mr.\ Donohue should treat
these as the bottom of Lux's negotiating range; any Satschel pushback
on them is a walk-away signal.

\textbf{1. Board composition --- 1 Lux director + 1 independent
director added to the existing 3-person Satschel board.} The Lux
director seat is for Lux Industries Inc.'s direct representative
(Mr.\ Dupont's nominee or Mr.\ Pahlavi's designate). The independent
director seat is a mutually agreed nominee with no prior relationship
to either side, serving as the ongoing governance anchor for the
SCLA terms. Effective on closing. No ownership-maintenance condition
attaches to either seat.

\textbf{2. 99-year exclusive USA-market license to Liquidity for
the LEL/LRL--PR portfolio in scope.} Conditional on Satschel
maintaining all SCLA covenants (royalty payments, anti-dilution,
governance terms, reciprocity grants). Termination of the SCLA for
material breach by Satschel terminates the exclusivity (license
reverts to non-exclusive). Outside the USA, Lux retains the right to
license the same IP to other parties.

\textbf{3. Anti-dilution: weighted-average broad-based, with full
ratchet on any down-round below the Lux issuance price.}

\textbf{4. Pro-rata participation right on all future Satschel
preferred financings.}

\textbf{5. Tag-along and co-sale rights on any founder transfer
exceeding 1\% of fully-diluted equity.}

\textbf{6. Information rights}: quarterly unaudited financials;
annual audited financials; copies of all board materials, minutes,
and consent resolutions; copies of all share-issuance and option-
grant resolutions; access to counsel and accountants under standard
NDA.

\textbf{7. Patent grant on all current and future Lux patents
covering the licensed IP.} Worldwide; perpetual; royalty-free under
the SCLA; expressly inclusive of any Lux patents not yet filed at
SCLA execution.

\textbf{8. Maintenance and support of upstream Lux OSS}: continuing
on the same shared-engineering basis as today, with no separate
maintenance fee.

\textbf{9. Reciprocity}: Lux-network L1/L2/L3/L4 chains that maintain
Lux consensus and Warp interop receive free non-exclusive license to
the Liquidity protocol (ATS/BD/TA). Lux chains may offer Liquidity-
style securities products without separate license.

\textbf{10. Engineering commitment from Satschel}: activate the Lux
Warp precompile (canonical address
\texttt{0x0200000000000000000000000000000000000005}) on every Liquid
chain at the next chain-upgrade window. New Liquid chains include
the \texttt{warpConfig} block in genesis.

\textbf{11. Litigation hold and document preservation}: Satschel
agrees to preserve all records, communications, and code review
artifacts relating to the licensing posture from 2014 to the SCLA
execution date.

\textbf{12. Independent investigation cooperation}: Satschel
cooperates fully with any independent investigation of officer
conduct undertaken by the Satschel disinterested directors or their
appointed counsel.

% ─────────────────────────────────────────────────────────────────
\section*{XII.\ Special Note Re: Mr.\ Zach Kelling}

\textbf{Mr.\ Kelling occupies a position requiring particular care
in Lux's litigation and settlement strategy.} The factual context
established in \S IV(F) bears directly here: Mr.\ Kelling was the
\emph{remediator} of the pre-Kelling Satschel technical and security
failures (the broken time-price priority matching; the \$40,000 ACH-
exploit theft) and the architect of the rebuilt platform that
subsequently supported the \$1.1B raise. He was not the author of
the conduct giving rise to this matter; he is, if anything, the
party whose work was misappropriated alongside Lux's. The further
facts:

\begin{itemize}[leftmargin=1.5em,itemsep=0.1em]
\item Mr.\ Kelling founded both Lux Industries Inc.\ and Satschel,
Inc., and authored substantial portions of the Lux IP at issue.
\item He is currently the Architect and Chief Cryptographer of
Satschel, but \emph{not} its CEO or CTO. He previously held those
titles in earlier corporate periods.
\item He is a \textbf{minority shareholder of Lux Industries Inc.}
\item He holds \textbf{zero equity in Satschel / Liquidity.io.}
\item He is not a director on either side.
\item He has been operationally displaced from Satschel's
commercial direction; the operational conduct giving rise to this
matter is attributable to other officers (Messrs.\ Choi, Trombley,
Church, and Sonsurkar) and not to him.
\end{itemize}

\textbf{Litigation-strategy implications:}

\begin{itemize}[leftmargin=1.5em,itemsep=0.1em]
\item Mr.\ Kelling should \emph{not} be named as a defendant in any
Lux suit. The conduct alleged here is not his.
\item Mr.\ Kelling is likely a material witness for Lux as to (a) the
authorship and ownership of the Lux IP, (b) the absence of any
executed commercial license between Lux and Satschel, and (c) the
operational displacement narrative explaining why the licensing
mismatch was not raised internally on his initiative.
\item Counsel should consider whether a litigation-cooperation
agreement, witness-protection arrangement, or formal indemnification
of Mr.\ Kelling (as a non-defendant witness) is appropriate to
ensure his continued availability as a witness without undue
personal exposure.
\item In any public communication regarding the matter, Lux should
take care not to attribute the misconduct to Mr.\ Kelling. The
narrative should be specific to the Satschel officers actually
responsible.
\item In the SCLA itself, Lux should consider providing for Mr.\
Kelling's protective treatment, including (a) a release of any
claims Lux might assert against him in his personal capacity, (b)
acknowledgement of his historical authorship contribution, and (c) a
ROFR or similar mechanism if his Satschel architectural role
generates IP that should flow back to Lux upstream.
\end{itemize}

\textbf{Practical effect}: Mr.\ Kelling can serve as a productive
intermediary in the Lux--Satschel negotiation, given his founder
status at both entities, his historical relationships, and his
absence of financial interest in Satschel that would compromise his
neutrality. Mr.\ Dupont should consider engaging Mr.\ Kelling
directly in informational meetings during the negotiation phase,
subject to outside counsel's oversight.

% ─────────────────────────────────────────────────────────────────
\section*{XIII.\ Procedural Recommendations}

\textbf{Decision chain:} Lux's decision flow on this matter is:

\begin{enumerate}[leftmargin=1.5em,itemsep=0.1em]
\item \textbf{Mr.\ Pahlavi (President \& CEO)} --- final approval
authority on any executed SCLA, any litigation filing, and any
strategic public communication.
\item \textbf{Mr.\ Dupont (Managing Partner, Strategic Transactions;
largest investor)} --- deal lead. Conducts the negotiation, selects
the settlement tier, and presents the executed term sheet to Mr.\
Pahlavi.
\item \textbf{Mr.\ Donohue (General Counsel)} --- legal lead.
Drafts the SCLA, supervises outside counsel, manages litigation
hold and document preservation, leads any pre-suit investigation,
and prepares the SCLA for execution by Mr.\ Pahlavi.
\item \textbf{Mr.\ Bin (CTO; this memorandum's author)} ---
operational and technical support. Provides factual development on
the IP portfolio, the licensing record, and the Satschel-side
operational context.
\end{enumerate}

\textbf{Recommended outside engagement:}

\begin{itemize}[leftmargin=1.5em,itemsep=0.1em]
\item \textbf{Outside litigation counsel}: engage a firm with
demonstrated DTSA + patent + Delaware-fiduciary-duty experience.
Recommended: Wilson Sonsini, Cooley, Fenwick \& West, or Quinn
Emanuel (for the litigation posture).
\item \textbf{Valuation firm}: engage an independent firm to opine
on the fairness of any equity-for-IP allocation in the SCLA.
Recommended: Houlihan Lokey, Duff \& Phelps (Kroll), or a specialist
boutique with digital-asset IP experience.
\item \textbf{Investigative firm}: engage a separate firm to
investigate Satschel officer conduct during the relevant period
(repository access, code review, internal communications, license
discussions). The investigation report informs both Lux's settlement
posture and any subsequent litigation. Recommended: Kroll
Investigations, FTI Consulting, or specialist tech-IP investigators.
\end{itemize}

\textbf{Recommended timeline:}

\begin{center}
\begin{tabular}{lp{10cm}}
\toprule
\textbf{Day} & \textbf{Action} \\
\midrule
0   & Mr.\ Pahlavi approves the memorandum and authorizes the deal team to proceed. \\
3   & Outside counsel engaged. \\
7   & Litigation hold notice transmitted to Satschel and named officers. \\
14  & Investigation firm begins document review. \\
21  & Term sheet delivered to Satschel for the recommended settlement tier. \\
30  & Satschel disinterested-director response received. \\
45  & Investigation report finalized; settlement terms tightened or expanded based on findings. \\
60  & Definitive SCLA executed. \\
75  & Closing: cash payment, share issuance, license activation. \\
\bottomrule
\end{tabular}
\end{center}

\textbf{If the timeline slips materially}, the deal team should
escalate to Mr.\ Pahlavi for guidance on whether to (a) extend
negotiation, (b) pivot to Path B leverage (a more aggressive
litigation-credible posture), or (c) file a protective complaint to
preserve statutory rights while continuing negotiation.

% ─────────────────────────────────────────────────────────────────
\section*{XIV.\ Prayer for Relief (Template)}

In the event Lux Industries Inc.\ proceeds to litigation, the prayer
for relief in the complaint will request:

\begin{enumerate}[leftmargin=1.5em,itemsep=0.1em]
\item Compensatory damages on each Count, in an amount to be proven
at trial, anticipated in the Small-to-Big Scenario range of \$80M to
\$590M.
\item Disgorgement of unjust enrichment derived from Satschel's
\$1.1 billion capital raise, in an amount attributable to the
misappropriated Lux IP.
\item Reasonable royalty under 35 U.S.C.\ \S 284 (patent) and 18
U.S.C.\ \S 1836(b)(3)(B) (DTSA).
\item Exemplary damages of 2$\times$ on the DTSA count (18 U.S.C.\
\S 1836(b)(3)(C)) for willful and malicious misappropriation.
\item Treble damages on the patent count (35 U.S.C.\ \S 284) for
willful infringement.
\item Statutory damages on the copyright count (17 U.S.C.\ \S 504(c))
at the elected maximum.
\item Attorneys' fees and costs under 18 U.S.C.\ \S 1836(b)(3)(D),
17 U.S.C.\ \S 505, 35 U.S.C.\ \S 285, and 15 U.S.C.\ \S 1117(a).
\item Preliminary and permanent injunctive relief enjoining Satschel
from further use of the misappropriated IP absent an executed
commercial license.
\item Pre-judgment and post-judgment interest.
\item Such other and further relief as the Court deems just and
proper.
\end{enumerate}

\textbf{Demand for jury trial} on all issues so triable.

% ─────────────────────────────────────────────────────────────────
\vspace{1em}
\hrule height 0.4pt
\vspace{0.5em}
\begin{center}
\textbf{\textsc{Schedules}}
\end{center}
\vspace{0.5em}

\section*{Schedule A --- Lux IP Inventory in Production at Satschel}

The full inventory is maintained at
\texttt{\textasciitilde/work/lux/legal/LIQUIDITY\_LICENSE\_AUDIT\_2026\_05\_25.md}.
Headline:

\begin{itemize}[leftmargin=1.5em,itemsep=0.1em]
\item BSD-3/Apache-2.0/MIT modules (no Lux action): approximately 25
modules.
\item LEL modules (commercial license required absent Authorized-
Network qualification): \texttt{aml}, \texttt{broker},
\texttt{captable}, \texttt{compliance}, \texttt{corona},
\texttt{crypto}, \texttt{fhe}, \texttt{hsm}, \texttt{maker},
\texttt{mpc}, \texttt{precompile} (and sub-precompiles),
\texttt{ringtail}, \texttt{threshold}, \texttt{transfer},
\texttt{treasury}, \texttt{warp}.
\item LRL--PR modules (research-only by default; commercial license
required): \texttt{vm}, \texttt{staking}, \texttt{sampler},
\texttt{dex} (Go bindings).
\item LGPL modules (compliance obligations independent of commercial
license): \texttt{evm}.
\item Modules without a LICENSE file (default copyright; no rights
granted absent affirmative license): \texttt{cex} and a small number
of related modules; counsel should treat as no-rights-granted pending
upstream remediation.
\end{itemize}

\section*{Schedule B --- Lux IP Pipeline (Not Yet in Production at
Satschel, Available Under SCLA)}

\begin{itemize}[leftmargin=1.5em,itemsep=0.1em]
\item AI inference and compute precompiles (LP-9080 family).
\item Advanced ZK rollup substrates.
\item Post-quantum hybrid signature schemes (X-Wing + ML-KEM; hybrid
ML-DSA + ECDSA).
\item Threshold MPC dynamic re-sharing protocols.
\item Confidential-execution attestation (TEE / SEV-SNP / TDX binding
to consensus).
\item White-label partner portal infrastructure (currently powering
MLC, VCC, TZERO, ReTokens).
\item Closed \texttt{\textasciitilde/work/luxcpp/*} (CUDA / Metal /
WebGPU / FPGA GPU accelerators, C++ DEX matching engine, AI compute
kernels) --- available under separate Private-Moat Addendum.
\end{itemize}

\section*{Schedule C --- Banking License Cost-to-Replicate (SFPB
Equivalent)}

\begin{center}
\begin{tabular}{lr}
\toprule
\textbf{License} & \textbf{Estimated Cost} \\
\midrule
US national or state banking charter & \$20--50M \\
Canadian Schedule II / partnership & \$10--20M \\
MAS Capital Markets Services + banking (Singapore) & \$5--10M \\
RBNZ deposit-taker (New Zealand) & \$5--15M \\
Sweden Finansinspektionen (EU passport) & \$5--10M \\
UK FCA EMI / Payment Institution & \$2--5M \\
Switzerland FINMA digital-asset banking & \$5--15M \\
Hong Kong SFC Type 1 / Type 7 + VATP & \$5--15M \\
\midrule
\textbf{Total direct regulatory + capital} & \textbf{\$57--140M} \\
Plus 5--7 years of timeline + execution risk & \\
\bottomrule
\end{tabular}
\end{center}

\section*{Schedule D --- Hanzo Inc.\ Portfolio Precedents}

\begin{itemize}[leftmargin=1.5em,itemsep=0.1em]
\item Damon Motorcycles --- electric / smart-motorcycle company.
\item Triller --- short-form video / social platform.
\item Bellabeat --- women's health wearables and AI-driven personal
health platform.
\item Unicorn Gold --- precious-metals tokenization / investment
platform.
\item Lux Network (Lux Industries Inc.) --- subject of this memorandum.
\item Zoo Labs Foundation / Zoo Network --- decentralized AI research
network.
\end{itemize}

Retained-parent-equity range across the portfolio: 10\%--40\% in
spawned operating companies. The 10\% ask in Satschel is at the
bottom of that internal range.

\section*{Schedule E --- Mandatory Terms (Lux's Floor)}

See \S XI, above.

\vfill
\hrule height 0.4pt
\vspace{0.5em}
\begin{small}
\textit{End of memorandum. This document is confidential, privileged,
and prepared in anticipation of litigation as well as for board
strategy purposes. The proposed settlement framework contained
herein is non-binding until executed as a definitive Settlement and
Cross-License Agreement. Distribution outside Lux Industries Inc.,
its Board of Directors, the named decision-chain officers (Messrs.\
Pahlavi, Dupont, Donohue, and Bin), and engaged outside counsel
requires the express written authorization of Mr.\ Pahlavi and Mr.\
Donohue.}
\end{small}

\end{document}
